\uuid{z9Rf}
\exo7id{3061}
\titre{exo7 3061}
\auteur{quercia}
\organisation{exo7}
\datecreate{2010-03-08}
\isIndication{false}
\isCorrection{false}
\chapitre{Propriétés de R}
\sousChapitre{Autre}
\cpge{mpsi-mp2i.structures-algebriques.c, mpsi-mp2i.nombres-reels-suites.a, pcsi.raisonnement-vocabulaire-ensembliste.c, ptsi.raisonnement-vocabulaire-ensembliste.c, mpsi-mp2i.nombres-reels-suites.g, bcpst1.nombres-reels, mp-mpi.topologie-espaces-normes.e, pc.espaces-vectoriels-normes.e, psi.espaces-vectoriels-normes.e, pt.fonctions-deux-trois-variables.b}
\module{Analyse}
\niveau{L1}
\difficulte{}

\contenu{
\texte{
Soit $f : \R \to \R$ un morphisme de corps.
}
\begin{enumerate}
    \item \question{Montrer que : $\forall\ x \in \Q,\ f(x) = x$.}
    \item \question{Montrer que $f$ est une application croissante.}
    \item \question{En d{\'e}duire que $f = \mathrm{id}_\R$.}
\end{enumerate}
}
